\section{Initial data and the main theorems}\label{sec:statements}

We use units $G=c=1$ and zero cosmological constant. The sign convention
for a spacetime realization is $(-,+,+,+)$ and
\[
 K(U,V)=\mathbf G(\nabla_U n,V),
\]
where $n$ is the future unit timelike normal. On an inner boundary, the
unit normal $\nu$ always points into the exterior, toward its designated
asymptotically Euclidean end. We put $H=\operatorname{div}_S\nu$ and
$\operatorname{tr}_S K=(g^{ij}-\nu^i\nu^j)K_{ij}$. Thus the future outward
null expansion is $\theta_+=H+\operatorname{tr}_S K$.

\begin{definition}[Charged exterior data]\label{def:data}
Let $\Omega$ be a connected oriented smooth three-manifold with nonempty
compact smooth boundary $S$, and let $(g,K,\EE,\BB)$ be smooth up to $S$.
Here $g$ is Riemannian, $K$ is a symmetric covariant two-tensor, and
$\EE,\BB$ are vector fields. We require completeness of the metric space
$(\Omega,g)$ with its boundary included. Outside a compact set there is
exactly one coordinate end, diffeomorphic to
$\{x\in\R^3:|x|>R_0\}$.

Throughout the proof we use geometric constraint densities:
\begin{align}
 2\mu&=R_g+\tau^2-|K|_g^2, &
 J_i&=\nabla^j(K_{ij}-\tau g_{ij}), &\tau&=\operatorname{tr}_gK,
 \label{eq:constraints}\\
 \mu_m&=\mu-(|\EE|_g^2+|\BB|_g^2), &
 J_m&=J-2(\EE\times\BB)^\flat.&&
 \label{eq:matter}
\end{align}
These four densities are $8\pi$ times their physical counterparts.
The cross product uses the orientation of $\Omega$. Assume
\begin{equation}\label{eq:charged-dec}
 \mu_m\ge |J_m|_g,\qquad
 \operatorname{div}_g\EE=\operatorname{div}_g\BB=0,
 \qquad \mu,|J|_g\in L^1(\Omega,dV_g).
\end{equation}
On the end, for some $q>1/2$ and $r=|x|$, assume
\begin{equation}\label{eq:weak-decay}
 g_{ij}-\delta_{ij}=O_2(r^{-q}),\qquad
 K_{ij}=O_1(r^{-1-q}),\qquad
 \EE^i,\BB^i=O_1(r^{-2}).
\end{equation}
The notation $O_j(r^{-a})$ bounds every coordinate derivative of order
$\ell\le j$ by $C_\ell r^{-a-\ell}$. These are bounds on otherwise
smooth data, not additional higher-derivative decay assumptions.

The following limits are assumed finite:
\begin{align}
 E&=\frac1{16\pi}\lim_{R\to\infty}
 \int_{r=R}(\partial_jg_{ij}-\partial_ig_{jj})n_\delta^i\,dA_\delta,
 \label{eq:adm-energy}\\
 P_i&=\frac1{8\pi}\lim_{R\to\infty}
 \int_{r=R}(K_{ij}-\tau g_{ij})n_\delta^j\,dA_\delta,
 \label{eq:adm-momentum}\\
 Q_E&=\frac1{4\pi}\lim_{R\to\infty}
 \int_{r=R}g(\EE,\nu_R)\,dA_g,
 &Q_B&=\frac1{4\pi}\lim_{R\to\infty}
 \int_{r=R}g(\BB,\nu_R)\,dA_g.
 \label{eq:charges}
\end{align}
Here $n_\delta,dA_\delta$ are Euclidean, whereas $\nu_R,dA_g$ are
computed in $g$. Put $Q=(Q_E^2+Q_B^2)^{1/2}$.
\end{definition}

No assumption is made about data behind $S$, and no extension across $S$
is required. In particular, the single-end assumption concerns this
exterior only. Boundary components and compact-interior topology are
unrestricted. There is no assumption of a global electromagnetic vector
potential. It is useful to retain instead the closed two-forms
\begin{equation}\label{eq:flux-forms}
 \alpha_1=\iota_{\EE}dV_g,\qquad
 \alpha_2=\iota_{\BB}dV_g.
\end{equation}
Unless stated otherwise, a change of metric keeps these forms fixed and
defines the new vector fields by contraction with the new volume form.
All form norms are exterior-algebra norms.

\begin{definition}[Full enclosing cut]\label{def:full-cut}
A full enclosing cut is $\Gamma=\partial D$, the entire intrinsic
manifold boundary of a connected smooth codimension-zero
submanifold-with-boundary $D\subset\Omega$ that is closed in $\Omega$,
has manifold interior in $\operatorname{int}\Omega$, and contains the
whole sufficiently distant part of the end. The boundary $\Gamma$ is
required to be compact, smooth, embedded and two-sided. It is oriented
toward the end. All components, including any portions coincident with
$S$, are counted. In particular $D=\Omega$ contributes $\Gamma=S$.
Define, in the original metric,
\begin{equation}\label{eq:enclosing-area}
 a_g(S)=\inf_\Gamma\Area_g(\Gamma),\qquad
 r_A=\sqrt{\frac{a_g(S)}{4\pi}}.
\end{equation}
\end{definition}

The closed flux forms have flux $4\pi Q_E$ and $4\pi Q_B$ across every
full cut by Stokes' Theorem. Opposite charges on different boundary
components are allowed. Section~\ref{sec:cuts} proves both positivity of
$a_g(S)$ and equivalence with the filled finite-perimeter convention used
in the proof; neither a positive infimum nor an attained smooth minimum
is assumed here.

\begin{theorem}[Charged spacetime Penrose inequality]\label{thm:numerical}
Let the data satisfy Definition~\ref{def:data}, and suppose
$\theta_+\le0$ on every component of $S$. Then $E>|P|$, the invariant
ADM mass $m=\sqrt{E^2-|P|^2}$ is positive, and
\begin{equation}\label{eq:main-inequality}
 E\ge\sqrt{|P|^2+Q^2},\qquad
 m\ge Q,\qquad
 r_A\le m+\sqrt{m^2-Q^2}.
\end{equation}
The same conclusion holds if, on each boundary component, one fixes
either the condition $H+\operatorname{tr}_S K\le0$ or the condition
$H-\operatorname{tr}_S K\le0$.
\end{theorem}

The upper-area formulation in \eqref{eq:main-inequality} is essential.
It is equivalent to $m\ge Q$ for $0<r_A\le Q$, and to
\begin{equation}\label{eq:polynomial-branch}
 m\ge\frac12\left(r_A+\frac{Q^2}{r_A}\right)
 \qquad\text{when }r_A>Q.
\end{equation}
There is no unrestricted polynomial claim below the charge threshold.
The theorem also proves that the possible null ADM endpoint cannot occur
with a nonempty compact boundary in this class. It requires no
connectedness, outermostness or area-minimization hypothesis for the
numerical inequality.

\begin{definition}[Connected horizon class]\label{def:rigidity-class}
Data in Theorem~\ref{thm:numerical} belong to the connected horizon class
if $S$ is connected, $\theta_+=0$ on $S$, every full cut has area at least
$A=\Area_g(S)>0$, and there is no compact smooth embedded two-sided full
enclosing surface entirely in $\operatorname{int}\Omega$ with
$\theta_+\le0$ everywhere. Such an excluded surface may be disconnected.
Thus $a_g(S)=A$. No strictly positive stability eigenvalue of this
marginally outer trapped surface (MOTS) is assumed.
\end{definition}

\begin{theorem}[Original-data equality]\label{thm:rigidity}
Assume the connected horizon class and $m>Q$. If
\begin{equation}\label{eq:main-equality}
 \sqrt{\frac{A}{4\pi}}=m+\sqrt{m^2-Q^2},
\end{equation}
then the entire original quadruple $(g,K,\EE,\BB)$, including its
boundary, is induced by a global smooth orientation-preserving spacelike
embedding into the regular future-horizon extension of one exterior of
dyonic Reissner--Nordstr\"om with parameters $(m,Q_E,Q_B)$.

The image lies in the domain of outer communication together with its
corresponding future horizon. The boundary is a full smooth horizon
cross-section or the bifurcation sphere, and the chosen end approaches
the corresponding spatial infinity. The embedding, its future unit
normal and the induced electromagnetic fields extend smoothly to $S$.
The non-electromagnetic matter densities vanish throughout the exterior.
\end{theorem}

For clarity, in static coordinates the target spacetime and two-form are
\begin{align}
 \mathbf G_{\rm RN}&=-F(r)\,dT^2+F(r)^{-1}dr^2+r^2\sigma_2,
 &F(r)&=1-\frac{2m}{r}+\frac{Q_E^2+Q_B^2}{r^2},
 \label{eq:rn-metric}\\
 \mathbf F_{\rm RN}&=\frac{Q_E}{r^2}\,dT\wedge dr
             +Q_B\,dA_{\sigma_2},
 &r&>r_+=m+\sqrt{m^2-Q^2}.
 \label{eq:rn-field}
\end{align}
Here $\sigma_2$ is the unit round metric with its oriented area form,
and $\operatorname{vol}_{\mathbf G}=r^2dT\wedge dr\wedge dA_{\sigma_2}$.
The Hodge star is characterized by
$\beta\wedge(*\gamma)=\langle\beta,\gamma\rangle_{\mathbf G}
\operatorname{vol}_{\mathbf G}$.
For the embedding $\iota$ and its future normal $n$, electromagnetic
matching means
\begin{equation}\label{eq:em-matching}
 \EE^\flat=\iota^*(\iota_n\mathbf F_{\rm RN}),\qquad
 \BB^\flat=\iota^*(\iota_n(*\mathbf F_{\rm RN})).
\end{equation}
The induced orientation is $\iota_n\operatorname{vol}_{\mathbf G}$.
On a static slice these fields are respectively $Q_E/r^2$ and $Q_B/r^2$
times the outward unit radial vector. On general slices neither radiality
nor pointwise parallelism of the original two fields is asserted.

\begin{theorem}[Converse and sharpness]\label{thm:converse}
Let $M>\sqrt{q_E^2+q_B^2}$, and put
$R_+=M+\sqrt{M^2-q_E^2-q_B^2}$. Every smooth spacelike exterior
hypersurface of dyonic Reissner--Nordstr\"om with these parameters whose
induced data satisfy Definitions~\ref{def:data} and
\ref{def:rigidity-class}, whose boundary is a full cross-section of the
corresponding future horizon or its bifurcation sphere, and whose actual
invariant ADM mass and outward charge fluxes are $M,q_E,q_B$, satisfies
\[
 a_g(S)=\Area_g(S)=4\pi R_+^2
\]
and attains equality in Theorem~\ref{thm:numerical}. For every such
parameter triple there are admissible static examples, examples with
$K\not\equiv0$, and examples with nonzero ADM momentum.
\end{theorem}

Mass and flux matching are explicit conditions in the converse;
the examples verify them by direct calculation. Theorem~\ref{thm:rigidity}
makes no extremal classification at $m=Q$ and no statement about data
behind the boundary.
